% !TeX program = xelatex
% UTF-8; compile twice with XeLaTeX. No external images or data files needed.
\documentclass[a4paper,landscape,10pt]{article}
\usepackage[margin=14mm,headheight=14pt,headsep=5mm,footskip=7mm]{geometry}
\usepackage{fontspec}
\setmainfont{DejaVu Sans}
\newfontfamily\ipafont{DejaVu Sans}
\usepackage[table]{xcolor}
\usepackage{array,longtable,booktabs,fancyhdr}
\definecolor{ink}{HTML}{18384D}
\definecolor{band}{HTML}{E7EFF4}
\pagestyle{fancy}\fancyhf{}
\fancyhead[L]{\small\color{ink}Parlez Vous QSO?! · Magyar / Hungarian}
\fancyhead[R]{\small Englische Phrase · Magyar fordítás · IPA}
\fancyfoot[L]{\scriptsize Vorlage: Hungarian v 1.0 · Credits: Dallos -- HA7PL}
\fancyfoot[R]{\small\thepage}
\setlength{\parindent}{0pt}
\setlength{\tabcolsep}{3mm}
\renewcommand{\arraystretch}{1.25}
\newcommand{\ipa}[1]{{\ipafont[#1]}}
\begin{document}
{\LARGE\bfseries\color{ink}Parlez Vous QSO?!}\hfill{\large Magyar / Hungarian}

\vspace{2mm}
{\small QSO-Sprachführer mit ungarischer IPA-Lautschrift. Die Lautschrift bezieht sich auf die ungarische Spalte.}

\vspace{2mm}
{\footnotesize \textbf{Leseschlüssel:} \ipa{ˈ} Wortbetonung, \ipa{ː} langer Vokal oder Konsonant;
\ipa{ɒ} ungarisches a, \ipa{aː} á, \ipa{ɛ} e, \ipa{eː} é; \ipa{ʃ} s, \ipa{s} sz, \ipa{ɲ} ny, \ipa{ɟ} gy, \ipa{c} ty.
Ungarische Wörter sind grundsätzlich auf der ersten Silbe betont. Schrägstriche trennen Alternativen; \ipa{…} steht für einen Platzhalter.}

\vspace{1mm}
\fontsize{9}{11.8}\selectfont
\begin{longtable}{@{}>{\raggedright\arraybackslash}p{13mm}>{\raggedright\arraybackslash}p{73mm}>{\raggedright\arraybackslash}p{75mm}>{\raggedright\arraybackslash}p{90mm}@{}}
\toprule
\textbf{Ref.} & \textbf{English phrase} & \textbf{Magyar} & \textbf{Lautschrift (IPA)}\\
\midrule
\endfirsthead
\toprule
\textbf{Ref.} & \textbf{English phrase} & \textbf{Magyar} & \textbf{Lautschrift (IPA)}\\
\midrule
\endhead
\midrule
\multicolumn{4}{r}{\scriptsize Fortsetzung auf der nächsten Seite}\\
\endfoot
\bottomrule
\endlastfoot
\rowcolor{band} 1.0 & \textbf{Opening Greetings} & \multicolumn{2}{l}{\textbf{Nyitó üdvözletek.}} \\*
1.1 & I am very pleased to meet you. & Nagyon örülök, hogy megismerhettem. & \ipa{ˈnɒɟon ˈøryløk hoɟ ˈmɛɡiʃmɛrhɛtːɛm} \\
1.2 & Thank you for your call & Köszönöm hívását. & \ipa{ˈkøsønøm ˈhiːvaːʃaːt} \\
1.3 & Good morning & Jó reggelt. & \ipa{ˈjoː ˈrɛɡːɛlt} \\
1.31 & Good day & Jó napot. & \ipa{ˈjoː ˈnɒpot} \\
1.32 & Good afternoon & Jó napot. \newline Jó délutánt & \ipa{ˈjoː ˈnɒpot / ˈjoː ˈdeːlutaːnt} \\
1.33 & Good evening & Jó estét. & \ipa{ˈjoː ˈɛʃteːt} \\
1.4 & I know only a little Hungarian but I enjoy QSOs in Hungarian. & Én csak keveset tudok magyarul, de örömmel QSO-zom magyarul & \ipa{ˈeːn ˈtʃɒk ˈkɛvɛʃɛt ˈtudok ˈmɒɟɒrul dɛ ˈørømːɛl ˈkuːɛʃoːzom ˈmɒɟɒrul} \\
\rowcolor{band} 2.0 & \textbf{Operator} & \multicolumn{2}{l}{\textbf{Operátor}} \\*
2.1 & My name is Peter. & Peter vagyok. \newline Az én nevem Péter & \ipa{ˈpiːtə ˈvɒɟok / ɒz ˈeːn ˈnɛvɛm ˈpeːtɛr} \\
2.2 & I have been licensed since 1960. & Én 1960 óta vagyok engedélyes. & \ipa{ˈeːn ˈɛzɛrkilɛntsːaːshɒtvɒn ˈoːtɒ ˈvɒɟok ˈɛŋɡɛdeːjɛʃ} \\
2.3 & I have been licensed for 43 years. & Én 43 éve vagyok engedélyes. & \ipa{ˈeːn ˈnɛɟvɛnhaːrom ˈeːvɛ ˈvɒɟok ˈɛŋɡɛdeːjɛʃ} \\
2.4 & I live in Sittingbourne. & Én Sittingbourne-ben élek. & \ipa{ˈeːn ˈsɪtɪŋbɔːn bɛn ˈeːlɛk} \\
2.5 & Sittingbourne is a village. & Sittingbourne egy falu. & \ipa{ˈsɪtɪŋbɔːn ɛɟː ˈfɒlu} \\
2.51 & Sittingbourne is a town. & Sittingbourne egy város. & \ipa{ˈsɪtɪŋbɔːn ɛɟː ˈvaːroʃ} \\
2.52 & Sittingbourne is a city. & Sittingbourne egy nagyváros. & \ipa{ˈsɪtɪŋbɔːn ɛɟː ˈnɒɟvaːroʃ} \\
2.6 & Sittingbourne is between London and Dover. & Sittingbourne London és Dover között található. & \ipa{ˈsɪtɪŋbɔːn ˈlondon eːʃ ˈdəʊvə ˈkøzøtː ˈtɒlaːlhɒtoː} \\
2.7 & I live in the South of England. & Én Anglia déli részén élek. & \ipa{ˈeːn ˈɒŋɡliɒ ˈdeːli ˈreːseːn ˈeːlɛk} \\
2.71 & I live in the North of England. & Én Anglia északi részén élek. & \ipa{ˈeːn ˈɒŋɡliɒ ˈeːsɒki ˈreːseːn ˈeːlɛk} \\
2.72 & I live in the East of England. & Én Anglia keleti részén élek. & \ipa{ˈeːn ˈɒŋɡliɒ ˈkɛlɛti ˈreːseːn ˈeːlɛk} \\
2.73 & I live in West of England. & Én Anglia nyugati részén élek. & \ipa{ˈeːn ˈɒŋɡliɒ ˈɲuɡɒti ˈreːseːn ˈeːlɛk} \\
2.74 & I live in South East England. & Én Dél-Kelet Angliában élek. & \ipa{ˈeːn ˈdeːlˌkɛlɛt ˈɒŋɡliaːbɒn ˈeːlɛk} \\
2.75 & I live in South West England. & Én Dél-Nyugat Angliában élek. & \ipa{ˈeːn ˈdeːlˌɲuɡɒt ˈɒŋɡliaːbɒn ˈeːlɛk} \\
2.76 & I live in North East England. & Én Észak-Kelet Angliában élek. & \ipa{ˈeːn ˈeːsɒkˌkɛlɛt ˈɒŋɡliaːbɒn ˈeːlɛk} \\
2.77 & I live in North West England. & Én Észak-Nyugat Angliában élek. & \ipa{ˈeːn ˈeːsɒkˌɲuɡɒt ˈɒŋɡliaːbɒn ˈeːlɛk} \\
2.78 & I live in central England. & Én Közép-Angliában élek. & \ipa{ˈeːn ˈkøzeːpˌɒŋɡliaːbɒn ˈeːlɛk} \\
2.8 & I live in the country. & Én falun élek. & \ipa{ˈeːn ˈfɒlun ˈeːlɛk} \\
2.81 & I live by the sea. & Én a tenger mellett élek. & \ipa{ˈeːn ɒ ˈtɛŋɡɛr ˈmɛlːɛtː ˈeːlɛk} \\
2.82 & I live near mountains. & Én a hegyek közelében élek. & \ipa{ˈeːn ɒ ˈhɛɟɛk ˈkøzɛleːbɛn ˈeːlɛk} \\
2.9 & Sittingbourne is 50Km from London. & Sittingbourne ötven kilométerre van Londontól. & \ipa{ˈsɪtɪŋbɔːn ˈøtvɛn ˈkilomeːtɛrːɛ vɒn ˈlondontoːl} \\
2.11 & I am retired. & Nyugdíjas vagyok. & \ipa{ˈɲuɡdiːjɒʃ ˈvɒɟok} \\
2.12 & My home call sign is M0ECS. & Az én otthoni hívójelem M0ECS. & \ipa{ɒz ˈeːn ˈotːhoni ˈhiːvoːjɛlɛm ɛm ˈnulːɒ eː tsɛː ɛʃ} \\
\rowcolor{band} 3.0 & \textbf{Station} & \multicolumn{2}{l}{\textbf{Az állomás}} \\*
3.1 & My radio is an Icom 706. & Az én rádióm egy Icom 706. & \ipa{ɒz ˈeːn ˈraːdioːm ɛɟː ˈaɪkɒm ˈheːtsːaːshɒt} \\
3.2 & My station is all homebrew (homemade). & Az én állomásom teljesen saját készítésű. & \ipa{ɒz ˈeːn ˈaːlːomaːʃom ˈtɛjːɛʃɛn ˈʃɒjaːt ˈkeːsiːteːʃyː} \\
3.3 & I am running 5 Watts. & Én 5 Wattal adok. & \ipa{ˈeːn ˈøt ˈvɒtːɒl ˈɒdok} \\
3.4 & My antenna is a long wire. & Az én antennám egy hosszú drót. & \ipa{ɒz ˈeːn ˈɒntɛnːaːm ɛɟː ˈhosːuː ˈdroːt} \\
3.41 & My antenna is a dipole. & Az én antennám egy dipól. & \ipa{ɒz ˈeːn ˈɒntɛnːaːm ɛɟː ˈdipoːl} \\
3.42 & My antenna is a G5RV. & Az én antennám egy G5RV. & \ipa{ɒz ˈeːn ˈɒntɛnːaːm ɛɟː ɡeː øt ɛr veː} \\
3.43 & My antenna is a doublet. & Az én antennám egy doublet. & \ipa{ɒz ˈeːn ˈɒntɛnːaːm ɛɟː ˈdʌblɪt} \\
3.44 & My antenna is a delta loop. & Az én antennám egy delta loop. & \ipa{ɒz ˈeːn ˈɒntɛnːaːm ɛɟː ˈdɛltɒ luːp} \\
3.45 & My antenna is a vertical. & Az én antennám egy vertikális antenna. & \ipa{ɒz ˈeːn ˈɒntɛnːaːm ɛɟː ˈvɛrtikaːliʃ ˈɒntɛnːɒ} \\
3.46 & My antenna is a log periodic. & Az én antennám egy log periodikus antenna. & \ipa{ɒz ˈeːn ˈɒntɛnːaːm ɛɟː ˈloɡ ˈpɛriodikuʃ ˈɒntɛnːɒ} \\
3.47 & My antenna is a Yagi beam. & Az én antennám egy Yagi beam. & \ipa{ɒz ˈeːn ˈɒntɛnːaːm ɛɟː ˈjɒɡi biːm} \\
3.48 & My antenna is a 5 element beam. & Az én antennám egy öt elemes beam. & \ipa{ɒz ˈeːn ˈɒntɛnːaːm ɛɟː ˈøt ˈɛlɛmɛʃ biːm} \\
\rowcolor{band} 4.0 & \textbf{Conditions} & \multicolumn{2}{l}{\textbf{Rádiózási feltételek}} \\*
4.1 & Thank you for my 599 report. & Köszönöm az 599 riportot. & \ipa{ˈkøsønøm ɒz ˈøt ˈkilɛnts ˈkilɛnts ˈriportot} \\
4.2 & The band conditions are good. & A sávkondíciók jók. & \ipa{ɒ ˈʃaːfkondiːtsioːk ˈjoːk} \\
4.3 & The band conditions are bad. & A sávkondíciók rosszak. & \ipa{ɒ ˈʃaːfkondiːtsioːk ˈrosːɒk} \\
4.4 & There is strong interference here. & Itt erős interferencia van. & \ipa{ˈitː ˈɛrøːʃ ˈintɛrfɛrɛntsijɒ vɒn} \\
4.5 & Can you change to the 80m band? & Át tud váltani a 80 méteres sávra? & \ipa{ˈaːt ˈtud ˈvaːltɒni ɒ ˈɲoltsvɒn ˈmeːtɛrɛʃ ˈʃaːvrɒ} \\
4.6 & Please change frequency to 3.558 MHz. & Kérem váltsunk frekvenciát 3.558 MHz. & \ipa{ˈkeːrɛm ˈvaːltʃuŋk ˈfrɛkvɛntsijaːt ˈhaːrom ˈɛɡeːs ˈøt ˈøt ˈɲolts ˈmɛɡɒhɛrts} \\
\rowcolor{band} 5.0 & \textbf{General Queries and Comments} & \multicolumn{2}{l}{\textbf{Általános kérdések és megjegyzések}} \\*
5.1 & What does <\#\#\#\#\#> mean please? & Mit jelent <\#\#\#\#\#> ez kérem? & \ipa{ˈmit ˈjɛlɛnt … ˈɛz ˈkeːrɛm} \\
5.2 & Please wait a moment, I need to check my Hungarian dictionary. & Kérem várjon egy pillanatot, meg kell néznem a magyar szótáramat & \ipa{ˈkeːrɛm ˈvaːrjon ɛɟː ˈpilːɒnɒtot ˈmɛɡ ˈkɛlː ˈneːznɛm ɒ ˈmɒɟɒr ˈsoːtaːrɒmɒt} \\
5.3 & How do you say, <\#\#\#\#\#> in Hungarian ? & Hogyan mondja <\#\#\#\#\#> magyarul? & \ipa{ˈhoɟɒn ˈmoɲɟːɒ … ˈmɒɟɒrul} \\
5.5 & My favourite modes of operation are CW / RTTY / PSK / MT63 / Olivia / Hellschreiber / Digital Modes / and SSB & Az én kedvenc üzemmódjaim: távíró, géptávíró / PSK / MT63 / Olívia,\newline Hellschreiber / digitális üzemmódok és az SSB (SSB=egyoldalsávos adás) & \ipa{ɒz ˈeːn ˈkɛdvɛnts ˈyzɛmːoːɟːɒim ˈtaːviːroː ˈɡeːptaːviːroː / peː ɛʃ kaː / ɛm teː ˈhɒtvɒnhaːrom / ˈoliːviɒ / ˈhɛlˌʃʁaɪbɐ / ˈdiɡitaːliʃ ˈyzɛmːoːdok eːʃ ɒz ɛʃ ɛʃ beː} \\
5.6 & My best DX is Australia. & A legjobb DX-em Ausztrália. & \ipa{ɒ ˈlɛɡjobː ˈdeːiksɛm ˈɒustraːliɒ} \\
5.7 & I enjoy working contests. & Szeretek versenyeken dolgozni. & \ipa{ˈsɛrɛtɛk ˈvɛrʃɛɲɛkɛn ˈdolɡozni} \\
5.8 & I like working QRP. & Kedvelem a QRP munkát. & \ipa{ˈkɛdvɛlɛm ɒ kuː ɛr peː ˈmuŋkaːt} \\
\rowcolor{band} 6.0 & \textbf{The Weather} & \multicolumn{2}{l}{\textbf{Az időjárás}} \\*
6.1 & It is sunny. & Napos. & \ipa{ˈnɒpoʃ} \\
6.2 & It is windy. & Szeles. & \ipa{ˈsɛlɛʃ} \\
6.3 & It is raining. & Esik az eső. & \ipa{ˈɛʃik ɒz ˈɛʃøː} \\
6.4 & It is snowing. & Havazik. & \ipa{ˈhɒvɒzik} \\
6.5 & It is very hot. & Nagyon forró. & \ipa{ˈnɒɟon ˈforːoː} \\
6.6 & It is very cold. & Nagyon hideg. & \ipa{ˈnɒɟon ˈhidɛɡ} \\
6.7 & The temperature is 8C. & A hőmérséklet 8 Celsius fok. & \ipa{ɒ ˈhøːmeːrʃeːklɛt ˈɲolts ˈtsɛlzijus ˈfok} \\
6.8 & It is overcast. & Borús. & \ipa{ˈboruːʃ} \\
6.9 & It is foggy. & Ködös. & \ipa{ˈkødøʃ} \\
6.11 & It has been raining all day. & Egész nap esett az eső. & \ipa{ˈɛɡeːs ˈnɒp ˈɛʃɛtː ɒz ˈɛʃøː} \\
6.12 & There is 20cm of snow here. & 20 cm hó van itt. & \ipa{ˈhuːs ˈtsɛntimeːtɛr ˈhoː vɒn ˈitː} \\
\rowcolor{band} 8.0 & \textbf{QSL} & \multicolumn{2}{l}{\textbf{QSL}} \\*
8.1 & My QSL card is sure via the bureau & Az én QSL lapom biztosan megy az irodán keresztül. & \ipa{ɒz ˈeːn kuː ɛʃ ɛl ˈlɒpom ˈbistoʃɒn ˈmɛɟ ɒz ˈirodaːn ˈkɛrɛstyl} \\
8.2 & QSL via eQSL & QSL lap e-QSL-en keresztül. & \ipa{kuː ɛʃ ɛl ˈlɒp iː kuː ɛʃ ɛlɛn ˈkɛrɛstyl} \\
8.3 & No eQSL please & Az e-QSL-t nem kérem. & \ipa{ɒz iː kuː ɛʃ ɛlt ˈnɛm ˈkeːrɛm} \\
8.4 & I will send you my card direct. & Én direkt fogom küldeni neked a lapomat. & \ipa{ˈeːn ˈdirɛkt ˈfoɡom ˈkyldɛni ˈnɛkɛd ɒ ˈlɒpomɒt} \\
8.5 & What is your address? & Mi a címed? & \ipa{ˈmi ɒ ˈtsiːmɛd} \\
8.6 & Please send your QSL card direct. & Kérlek, küldd a QSL lapodat direkt. & \ipa{ˈkeːrlɛk ˈkyldː ɒ kuː ɛʃ ɛl ˈlɒpodɒt ˈdirɛkt} \\
8.7 & My address is... & Az én címem… & \ipa{ɒz ˈeːn ˈtsiːmɛm …} \\
8.8 & My address is shown on QSL.com / EHam.net. & A címem látható a QSL.com / EHam.net. & \ipa{ɒ ˈtsiːmɛm ˈlaːthɒtoː ɒ kuː ɛʃ ɛl pont kom / iː hæm pont nɛt} \\
8.9 & Sorry I do not send QSL cards. & Sajnos én nem küldök QSL lapokat. & \ipa{ˈʃɒjnoʃ ˈeːn ˈnɛm ˈkyldøk kuː ɛʃ ɛl ˈlɒpokɒt} \\
8.11 & Please send QSL card via G1JQH. & Kérem, küldje a QSL lapot G1JQH-n keresztül. & \ipa{ˈkeːrɛm ˈkylɟːɛ ɒ kuː ɛʃ ɛl ˈlɒpot ɡeː ɛɟː jeː kuː haːn ˈkɛrɛstyl} \\
8.12 & My WACRAL / GQRP / DXCC / FISTS number is 1234. & Az én WACRAL / GQRP / DXCC / FISTS számom 1234. & \ipa{ɒz ˈeːn ˈvɒkrɒl / ɡeː kuː ɛr peː / deː iks tsɛː tsɛː / fɪsts ˈsaːmom ˈɛzɛrkeːtsːaːzˌhɒrmintsneːɟ} \\
\rowcolor{band} 9.0 & \textbf{Final Greetings} & \multicolumn{2}{l}{\textbf{Búcsúzó üdvözletek}} \\*
9.1 & Thank you. & Köszönöm neked. & \ipa{ˈkøsønøm ˈnɛkɛd} \\
9.2 & Thank you for a very enjoyable contact. & Köszönöm neked a nagyon kellemes összeköttetést. & \ipa{ˈkøsønøm ˈnɛkɛd ɒ ˈnɒɟon ˈkɛlːɛmɛʃ ˈøsːɛkøtːɛteːʃt} \\
9.3 & I enjoyed our chat. & Élveztem a beszélgetésünket. & \ipa{ˈeːlvɛztɛm ɒ ˈbɛseːlɡɛteːʃyŋkɛt} \\
9.4 & I enjoyed our QSO in Hungarian. & Élveztem a magyar nyelvű KUESÓNKAT. (QSO=összeköttetés) & \ipa{ˈeːlvɛztɛm ɒ ˈmɒɟɒr ˈɲɛlvyː ˈkuːɛʃoːŋkɒt} \\
9.5 & It has been great fun. & Nagy öröm volt számomra. & \ipa{ˈnɒɟ ˈørøm ˈvolt ˈsaːmomrɒ} \\
9.6 & If you would like a regular contact with me, you can contact me via email. & Ha rendszeres összeköttetést szeretnél velem, kapcsolatba léphetsz velem emailen keresztül. & \ipa{hɒ ˈrɛntsɛrɛʃ ˈøsːɛkøtːɛteːʃt ˈsɛrɛtneːl ˈvɛlɛm ˈkɒptʃolɒdbɒ ˈleːphɛts ˈvɛlɛm ˈiːmeːlɛn ˈkɛrɛstyl} \\
9.7 & My email address is... & Az én email címem… & \ipa{ɒz ˈeːn ˈiːmeːl ˈtsiːmɛm …} \\
9.8 & My email address is shown on QSL.com / EHam.net & Az email címem látható a QSL.com / EHam.net. & \ipa{ɒz ˈiːmeːl ˈtsiːmɛm ˈlaːthɒtoː ɒ kuː ɛʃ ɛl pont kom / iː hæm pont nɛt} \\
9.9 & I look forward to meeting you again. & Várom a következő találkozást veled. & \ipa{ˈvaːrom ɒ ˈkøvɛtkɛzøː ˈtɒlaːlkozaːʃt ˈvɛlɛd} \\
9.11 & My best wishes to you and your family. & Szívélyes üdvözletem neked és családodnak. & \ipa{ˈsiːveːjɛʃ ˈydvøzlɛtɛm ˈnɛkɛd eːʃ ˈtʃɒlaːdodnɒk} \\
9.12 & Best wishes to your wife. & Szívélyes üdvözletem a feleségednek. & \ipa{ˈsiːveːjɛʃ ˈydvøzlɛtɛm ɒ ˈfɛlɛʃeːɡɛdnɛk} \\
9.121 & Best wishes to your husband. & Szívélyes üdvözletem a férjednek. & \ipa{ˈsiːveːjɛʃ ˈydvøzlɛtɛm ɒ ˈfeːrjɛdnɛk} \\
9.13 & Good luck. & Jó szerencsét. & \ipa{ˈjoː ˈsɛrɛntʃeːt} \\
9.14 & Good health. & Jó egészséget. & \ipa{ˈjoː ˈɛɡeːʃːeːɡɛt} \\
9.15 & Goodbye & Viszontlátásra.\newline Isten veled. & \ipa{ˈvisontlaːtaːʃrɒ / ˈiʃtɛn ˈvɛlɛd} \\
9.16 & I wish you an enjoyable afternoon. & Kellemes délutánt kívánok neked. & \ipa{ˈkɛlːɛmɛʃ ˈdeːlutaːnt ˈkiːvaːnok ˈnɛkɛd} \\
9.161 & I wish you an enjoyable evening. & Kellemes estét kívánok neked. & \ipa{ˈkɛlːɛmɛʃ ˈɛʃteːt ˈkiːvaːnok ˈnɛkɛd} \\
9.162 & I wish you an enjoyable weekend. & Kellemes hétvégét kívánok neked. & \ipa{ˈkɛlːɛmɛʃ ˈheːtveːɡeːt ˈkiːvaːnok ˈnɛkɛd} \\
9.17 & Enjoy your holiday. & Élvezd a szabadságodat. & \ipa{ˈeːlvɛzd ɒ ˈsɒbɒtʃːaːɡodɒt} \\
9.18 & Happy New Year. & Boldog új évet. & \ipa{ˈboldoɡ ˈuːj ˈeːvɛt} \\
\rowcolor{band} 10.0 & \textbf{Country} & \multicolumn{2}{l}{\textbf{Országok}} \\*
10.1 & Albania & Albánia & \ipa{ˈɒlbaːniɒ} \\
10.2 & Austria & Ausztria & \ipa{ˈɒustriɒ} \\
10.3 & Belgium & Belgium & \ipa{ˈbɛlɡijum} \\
10.4 & Belorussia & Belorusszia & \ipa{ˈbɛlorusːiɒ} \\
10.5 & Bosnia \& Herzegovina & Bosznia és Hercegovina & \ipa{ˈbosniɒ eːʃ ˈhɛrtsɛɡovinɒ} \\
10.6 & Bulgaria & Bulgária & \ipa{ˈbulɡaːriɒ} \\
10.7 & Croatia & Horvátország & \ipa{ˈhorvaːtorsaːɡ} \\
10.8 & Czech Republic & Cseh Köztársaság & \ipa{ˈtʃɛ ˈkøztaːrʃɒʃaːɡ} \\
10.9 & Denmark & Dánia & \ipa{ˈdaːniɒ} \\
10.11 & Eire & Írország & \ipa{ˈiːrorsaːɡ} \\
10.12 & England & Anglia & \ipa{ˈɒŋɡliɒ} \\
10.13 & Estonia & Észtország & \ipa{ˈeːstorsaːɡ} \\
10.14 & Finland & Finnország & \ipa{ˈfinːorsaːɡ} \\
10.15 & France & Franciaország & \ipa{ˈfrɒntsijɒorsaːɡ} \\
10.16 & Germany & Németország & \ipa{ˈneːmɛtorsaːɡ} \\
10.17 & Gibraltar & Gibraltár & \ipa{ˈɡibrɒltaːr} \\
10.18 & Greece & Görögország & \ipa{ˈɡørøɡorsaːɡ} \\
10.19 & Holland & Hollandia & \ipa{ˈholːɒndiɒ} \\
10.21 & Hungary & Magyarország & \ipa{ˈmɒɟɒrorsaːɡ} \\
10.22 & Iceland & Izland & \ipa{ˈizlɒnd} \\
10.23 & Northern Ireland & Észak Írország & \ipa{ˈeːsɒk ˈiːrorsaːɡ} \\
10.24 & Isle of Man & Man sziget & \ipa{ˈmɛn ˈsiɡɛt} \\
10.25 & Italy & Olaszország & \ipa{ˈolɒsorsaːɡ} \\
10.26 & Latvia & Lettország & \ipa{ˈlɛtːorsaːɡ} \\
10.27 & Liechtenstein & Liechtenstein & \ipa{ˈlɪçtn̩ʃtaɪn} \\
10.28 & Lithuania & Litvánia & \ipa{ˈlitvaːniɒ} \\
10.29 & Luxembourg & Luxemburg & \ipa{ˈluksɛmburɡ} \\
10.31 & Macedonia & Macedónia & \ipa{ˈmɒtsɛdoːniɒ} \\
10.32 & Moldavia & Moldávia & \ipa{ˈmoldaːviɒ} \\
10.33 & Netherlands & Hollandia & \ipa{ˈholːɒndiɒ} \\
10.34 & Norway & Norvégia & \ipa{ˈnorveːɡiɒ} \\
10.35 & Poland & Lengyelország & \ipa{ˈlɛɲɟɛlorsaːɡ} \\
10.36 & Portugal & Portugália & \ipa{ˈportuɡaːliɒ} \\
10.37 & Romania & Románia & \ipa{ˈromaːniɒ} \\
10.38 & Russian Federation & Orosz Föderáció & \ipa{ˈoros ˈfødɛraːtsioː} \\
10.39 & Scotland & Skócia & \ipa{ˈʃkoːtsiɒ} \\
10.41 & Slovakia & Szlovákia & \ipa{ˈslovaːkiɒ} \\
10.42 & Slovenia & Szlovénia & \ipa{ˈsloveːniɒ} \\
10.43 & Spain & Spanyolország & \ipa{ˈʃpɒɲolorsaːɡ} \\
10.44 & Svalbard & Szvalbard & \ipa{ˈsvɒlbɒrd} \\
10.45 & Sweden & Svédország & \ipa{ˈʃveːdorsaːɡ} \\
10.46 & Switzerland & Svájc & \ipa{ˈʃvaːjts} \\
10.47 & Turkey & Törökország & \ipa{ˈtørøkorsaːɡ} \\
10.48 & Ukraine & Ukrajna & \ipa{ˈukrɒjnɒ} \\
10.49 & Wales & Wales & \ipa{ˈvɛls} \\
\end{longtable}
\normalsize
\section*{Hinweise zur Übertragung und Aussprache}
\fontsize{8.5}{10}\selectfont\raggedright
\textbf{Vorlage und Umfang.} Grundlage ist „Hungarian v 1.0.doc“; Credit: Dallos, HA7PL.
Alle 142 Phrasen und Ländernamen mit ihren Alternativen sind enthalten. Nummernsprünge und der fehlende Abschnitt 7 bleiben erhalten.

\smallskip
\textbf{Korrekturen.} Referenz „132“ wurde zu „1.32“. „belváros“ (Innenstadt) in 2.52 wurde gemäß „city“ zu
„nagyváros“ (Großstadt), „angliában“ zu „Angliában“ (2.75), „dipol“ zu „dipól“ (3.41).
Die Aufforderungen in 8.6 und 8.11 lauten nun „Kérlek, küldd“ und „Kérem, küldje“.
„Liktenstein“ wurde zu „Liechtenstein“ (10.27), „Littvánia“ zu „Litvánia“ (10.28), „Welsz“ zu „Wales“ (10.49)
und „Man Sziget“ zu „Man sziget“ (10.24). Und-Zeichen zwischen Alternativsätzen wurden durch Zeilenumbrüche ersetzt;
Mehrfachleerzeichen wurden vereinheitlicht.

\smallskip
\textbf{Quellenbesonderheiten.} Die ungewöhnliche Wortstellung in 5.1, der fehlende Anschluss an die
Frequenzangabe in 4.6 und „esett“ (es regnete) in 6.11 bleiben erhalten. Ländernamen und Webadressen wurden nicht modernisiert; Holland und Netherlands stehen beide als „Hollandia“.
Die Klammererklärungen in 5.5 und 9.4 erläutern Bedeutungen und werden in der IPA-Spalte nicht mitgesprochen.

\smallskip
\textbf{Aussprache und Betonung.} Die IPA-Spalte ist eine breite Aussprachehilfe für Standardungarisch.
Ungarische Wörter sind grundsätzlich auf der ersten Silbe betont; Artikel und einige Funktionswörter
stehen ohne eigenes Betonungszeichen. \ipa{ˌ} zeigt ausgewählte Nebenbetonungen in Zusammensetzungen.
Fremdnamen können anders betont sein; Satzmelodie ist nicht kodiert.

\smallskip
\textbf{Vokale.} \ipa{ː} markiert bedeutungsunterscheidende Länge, auch bei Konsonanten.
a/á stehen als \ipa{ɒ/aː}, e/é als \ipa{ɛ/eː}; diese Paare unterscheiden sich auch in der Lautqualität.
ö/ő = \ipa{ø/øː}, ü/ű = \ipa{y/yː}, i/í = \ipa{i/iː}, o/ó = \ipa{o/oː}, u/ú = \ipa{u/uː}.
Unbetonte Vokale werden nicht generell zu Schwa reduziert.

\smallskip
\textbf{Konsonanten.} s = \ipa{ʃ}, sz = \ipa{s}, zs = \ipa{ʒ}, cs = \ipa{tʃ}, c = \ipa{ts};
gy \ipa{ɟ}, ty \ipa{c} und ny \ipa{ɲ} sind palatale Laute. ly wird als \ipa{j} gesprochen.
\ipa{ŋ} steht für n vor k/g. Ausgewählte Angleichungen sind berücksichtigt, etwa „mondja“ mit \ipa{ɟː},
„egészséget“ mit \ipa{ʃː} und „szabadságodat“ mit \ipa{tʃː}. Das h in „Cseh“ ist stumm dargestellt.
Weitere Angleichungen über Wortgrenzen, Kürzungen langer Konsonanten in Gruppen und feine h-Varianten sind nicht vollständig abgebildet.

\smallskip
\textbf{Namen und Kürzel.} Sittingbourne, Dover und Peter (erste Variante von 2.1) werden englisch gelesen,
Péter (zweite Variante) ungarisch. Icom, Doublet, Loop und Beam verwenden englische, Hellschreiber und
Liechtenstein deutsche Varianten. Andere Namen sind überwiegend ungarisch, auch Wales („velsz“).
Kürzel und Rufzeichen verwenden ungarische Buchstabennamen: Q „kú“, S „es“, C „cé“, J „jé“, H „há“.
WACRAL und FISTS stehen als Wörter. QSO-zom und KUESÓNKAT entsprechen QSO mit ungarischen Endungen.
Das e in e-QSL/eHam ist \ipa{iː}; email wird als „ímél“, Webpunkte als „pont“, com und net als Wörter gelesen.

\smallskip
\textbf{Zahlen und Platzhalter.} 1960 = „ezerkilencszázhatvan“, 43 = „negyvenhárom“, 706 = „hétszázhat“,
MT63 mit „hatvanhárom“, 599 = „öt kilenc kilenc“, 1234 = „ezerkétszáz-harmincnégy“.
3.558 MHz wird als „három egész öt öt nyolc megahertz“ gelesen; 80 m und 20 cm mit „nyolcvan méter“ und „húsz centiméter“.
\texttt{<\#\#\#\#\#>} und offene Adressangaben erscheinen als Auslassung.

\smallskip
\textbf{Fachhinweis.} Tamás Szende: „Hungarian“, Journal of the International Phonetic Association
24(2), 1994, S. 91--94, DOI 10.1017/S0025100300005090 (bibliografischer Hinweis zur IPA-Beschreibung).
Namens- und Kürzelvarianten sind möglich.
\end{document}
